\documentclass[xcolor=dvipsnames,t,aspectratio=169]{beamer}

\usecolortheme{rose}
\usecolortheme{dolphin}
\usetheme{Boadilla}

\input{../../templates/slides/imports}
\input{../../templates/slides/settings}
\input{../../templates/slides/commands}

\titlegraphic{
    \includegraphics[scale = 0.5]{../../templates/slides/logo}
}

\logo{
\begin{tikzpicture}[overlay,remember picture]
\node[xshift=-.75cm, yshift=-1cm] at (current page.30){
    \includegraphics[width=0.1\textwidth]{../../templates/slides/logo}};
\end{tikzpicture}
}
\usetikzlibrary{positioning}

\newcommand{\highlight}[1]{{\color{nes_dark_orange} #1}}

\title{Clusterização, Busca Textual e Ranking}

\author{Eduardo Adame}

\date{{\color{nes_dark_purple}\textbf{Ciência de Dados}\\[0.5em] 22 de maio de 2026 }}

\begin{document}

\frame[plain]{\titlepage}
\setcounter{framenumber}{0}


% ============================================================
\section{Clusterização}
% ============================================================

\begin{frame}[c]{Aprendizado não supervisionado}

\begin{columns}
\column{0.6\textwidth}
\textbf{\color{nes_dark_purple}O que é clusterização?}
\begin{itemize}
    \item Agrupar pontos \highlight{sem rótulos} — o algoritmo descobre a estrutura
    \item Objetivo: maximizar similaridade \emph{dentro} do grupo e minimizar \emph{entre} grupos
    \item Aplicações: segmentação de clientes, detecção de anomalias, \highlight{churn}, compressão de dados, pré-processamento para classificação
\end{itemize}
\textbf{\color{nes_dark_purple}Famílias de algoritmos}
\begin{itemize}
    \item \highlight{Centróides}: K-Means, K-Medoids
    \item \highlight{Densidade}: DBSCAN, OPTICS, HDBSCAN
    \item \highlight{Hierárquicos}: aglomerativo, divisivo
    \item \highlight{Probabilísticos}: GMM, LDA
\end{itemize}
\column{0.4\textwidth}
\vspace*{.4cm}


\textbf{\color{nes_dark_purple}Churn — um caso de uso real}
\begin{itemize}
    \item Clientes com comportamento similar ao de quem cancelou são agrupados
    \item Cluster de \highlight{alto risco}: baixa frequência, tickets altos, sem renovação recente
    \item Cluster de \highlight{baixo risco}: uso crescente, suporte esporádico
    \item Permite agir \emph{antes} do cancelamento — retenção proativa
\end{itemize}
\end{columns}


\end{frame}


\begin{frame}[c]{K-Means — o algoritmo}

\begin{figure}
    \includegraphics[height=.78\textheight]{kmeans_steps.png}
\end{figure}

\end{frame}


\begin{frame}[c, fragile]{K-Means na prática}

\begin{code}[scikit-learn: K-Means e escolha de k]{python}
from sklearn.cluster import KMeans
from sklearn.preprocessing import StandardScaler
import matplotlib.pyplot as plt
# Sempre escalar antes de clusterizar
X_std = StandardScaler().fit_transform(X)
# Método do cotovelo — escolher k
inercias = []
for k in range(1, 11):
    km = KMeans(n_clusters=k, random_state=42, n_init=10)
    km.fit(X_std)
    inercias.append(km.inertia_)
plt.plot(range(1, 11), inercias, marker="o")
plt.xlabel("k"); plt.ylabel("Inércia"); plt.title("Método do Cotovelo")
\end{code}

\begin{itemize}
    \item \texttt{inertia\_} = soma das distâncias quadráticas ao centróide
    \item \highlight{Método do cotovelo}: escolha o $k$ onde a queda desacelera
    \item \texttt{n\_init=10} roda 10 inicializações aleatórias — usa a melhor
\end{itemize}

\end{frame}


\begin{frame}[c, fragile]{K-Means — avaliação e visualização}

\begin{code}[Silhouette score e redução dimensional]{python}
from sklearn.metrics import silhouette_score
from sklearn.decomposition import PCA
km = KMeans(n_clusters=3, random_state=42, n_init=10)
labels = km.fit_predict(X_std)
# Silhouette: [-1, 1] — quanto maior, mais coeso e separado
score = silhouette_score(X_std, labels)
print(f"Silhouette: {score:.3f}")
# Visualizar em 2D com PCA
pca = PCA(n_components=2)
X2d = pca.fit_transform(X_std)
plt.scatter(X2d[:, 0], X2d[:, 1], c=labels, cmap="tab10", alpha=0.6)
plt.scatter(pca.transform(km.cluster_centers_)[:, 0],
            pca.transform(km.cluster_centers_)[:, 1],
            marker="X", s=200, c="black", label="Centróides")
plt.legend()
\end{code}
\end{frame}


\begin{frame}[c]{DBSCAN — clusterização por densidade}

\begin{figure}
    \includegraphics[height=.78\textheight]{kmeans_vs_dbscan.jpg}
\end{figure}

\end{frame}


\begin{frame}[c, fragile]{DBSCAN na prática}

\begin{code}[scikit-learn: DBSCAN e HDBSCAN]{python}
from sklearn.cluster import DBSCAN
# eps: raio da vizinhança; min_samples: mínimo para ser núcleo
db = DBSCAN(eps=0.5, min_samples=5)
labels = db.fit_predict(X_std)
n_clusters = len(set(labels)) - (1 if -1 in labels else 0)
n_ruido    = (labels == -1).sum()
print(f"{n_clusters} clusters, {n_ruido} pontos de ruído")
# HDBSCAN — versão hierárquica, não precisa de eps
# pip install hdbscan   (ou sklearn >= 1.3)
from sklearn.cluster import HDBSCAN
hdb = HDBSCAN(min_cluster_size=10)
labels_h = hdb.fit_predict(X_std)
\end{code}

\begin{itemize}
    \item \texttt{-1} = ruído — ponto que não pertence a nenhum cluster
    \item \highlight{DBSCAN} encontra clusters de \emph{forma arbitrária} e detecta outliers
    \item \highlight{HDBSCAN} dispensa $\varepsilon$ — mais robusto para dados reais
\end{itemize}
\end{frame}



% ============================================================
\section{Busca textual e TF-IDF}
% ============================================================

\begin{frame}[c]{Do texto bruto ao vetor}

\begin{columns}
\column{0.5\textwidth}
\textbf{\color{nes_dark_purple}Pipeline de pré-processamento}
\begin{enumerate}
    \item \highlight{Tokenização} — dividir em tokens (palavras, subpalavras)
    \item \highlight{Normalização} — minúsculas, remoção de acentos
    \item \highlight{Stopwords} — remover palavras muito frequentes e pouco informativas
    \item \highlight{Stemming / Lematização} — reduzir à raiz morfológica
    \item \highlight{Vetorização} — converter tokens em números
\end{enumerate}

\column{0.5\textwidth}
\textbf{\color{nes_dark_purple}Representações vetoriais}
\begin{itemize}
    \item \highlight{Bag of Words (BoW)}: contagem de ocorrências — simples, sem contexto
    \item \highlight{TF-IDF}: pondera pela raridade no corpus — muito mais discriminativo
    \item \highlight{Word2Vec / GloVe}: vetores densos com semântica
    \item \highlight{Sentence Transformers}: embeddings de sentenças inteiras
\end{itemize}

% \vspace{0.4em}
Para busca e ranking: TF-IDF. Para semântica e similaridade: embeddings.
% \end{display}
\end{columns}

\end{frame}


\begin{frame}[c]{TF-IDF — a intuição}

\begin{figure}
    \includegraphics[height=.425\textheight]{tfidf_intuition.png}
\end{figure}

\end{frame}


\begin{frame}[c, fragile]{TF-IDF — implementação}

\begin{code}[scikit-learn: vetorização e busca por cosseno]{python}
from sklearn.feature_extraction.text import TfidfVectorizer
from sklearn.metrics.pairwise import cosine_similarity
corpus = ["gato comeu peixe dormiu",
          "cachorro correu parque latiu",
          "gato cachorro são animais domésticos"]
vec = TfidfVectorizer(max_features=1000, sublinear_tf=True)
X   = vec.fit_transform(corpus)      # matriz esparsa (docs × termos)
# Busca: encontrar documento mais similar à query
query  = "animal de estimação"
q_vec  = vec.transform([query])
scores = cosine_similarity(q_vec, X).flatten()
for i, score in sorted(enumerate(scores), key=lambda x: -x[1]):
    print(f"[{score:.3f}] {corpus[i]}")
\end{code}

\begin{itemize}
    \item \highlight{Similaridade por cosseno}: independe do comprimento do documento
    \item \texttt{sublinear\_tf=True}: usa $1 + \log(\text{tf})$ — suaviza termos muito frequentes
\end{itemize}

\end{frame}


\begin{frame}[c, fragile]{Tokenização e busca com banco de dados}

\begin{code}[SQLite + TF-IDF: buscador persistente]{python}
# 1. Armazenar documentos no banco
conn = sqlite3.connect("buscador.db")
conn.execute("CREATE TABLE IF NOT EXISTS docs "
             "(id INTEGER PRIMARY KEY, titulo TEXT, conteudo TEXT)")
conn.executemany("INSERT INTO docs (titulo, conteudo) VALUES (?,?)", docs)
conn.commit()
# 2. Treinar vetorizador e serializar
docs_texto = [row[0] for row in conn.execute("SELECT conteudo FROM docs")]
vec = TfidfVectorizer(max_features=5000, sublinear_tf=True)
X   = vec.fit_transform(docs_texto)
# 3. Buscar
def buscar(query, top_k=5):
    vec_q = vec.transform([query])
    scores = cosine_similarity(vec_q, X).flatten()
    top = scores.argsort()[::-1][:top_k]
    return conn.execute(
        f"SELECT titulo FROM docs WHERE id IN ({','.join('?'*top_k)})",
        (top + 1).tolist()).fetchall()
\end{code}

\end{frame}


% ============================================================
\section{PageRank}
% ============================================================

\begin{frame}[c]{A história do PageRank}

\begin{columns}
\column{0.5\textwidth}
\textbf{\color{nes_dark_purple}O problema de 1998}
\begin{itemize}
    \item Alta Visa e outros buscadores ranqueavam por \highlight{correspondência de palavras}
    \item Fácil de manipular: repetir a palavra-chave mil vezes
    \item Sergey Brin e Larry Page (Stanford) perguntaram: \highlight{``e se um link fosse um voto?''}
    \item Insight: um voto de uma página importante vale mais que de uma irrelevante
\end{itemize}

\vspace{0.4em}
\textbf{\color{nes_dark_purple}A metáfora do internauta aleatório}
\begin{itemize}
    \item Imagine um usuário seguindo links \highlight{aleatoriamente}
    \item Com probabilidade $\alpha$ clica num link qualquer da página
    \item Com probabilidade $1 - \alpha$ \highlight{teleporta} para uma página aleatória
    \item PageRank = fração do tempo que o internauta passa em cada página
\end{itemize}

\column{0.5\textwidth}
\textbf{\color{nes_dark_purple}Definição formal}

Seja $N$ o número de páginas e $L(j)$ o conjunto de páginas que apontam para $j$:

$$\text{PR}(j) = \frac{1-\alpha}{N} + \alpha \sum_{i \in L(j)} \frac{\text{PR}(i)}{|C(i)|}$$

onde $|C(i)|$ é o número de links saindo de $i$ e $\alpha \approx 0{,}85$.

\vspace{0.5em}
\begin{display}[Convergência]
A iteração converge pois a matriz de transição é estocástica e irredutível graças ao fator de teleporte.
\end{display}
\end{columns}

\end{frame}


\begin{frame}[c]{PageRank — o internauta aleatório}

\begin{figure}
    \includegraphics[height=.78\textheight]{pagerank_graph.jpg}
\end{figure}

\end{frame}


\begin{frame}[c, fragile]{PageRank — implementação iterativa}

\begin{code}[NetworkX: PageRank em grafo da Wikipedia]{python}
import networkx as nx
import requests
from bs4 import BeautifulSoup
def links_da_pagina(titulo):
    url = f"https://pt.wikipedia.org/wiki/{titulo}"
    soup = BeautifulSoup(requests.get(url,
           headers={"User-Agent": "Mozilla/5.0"}).text, "lxml")
    return [a["href"].replace("/wiki/", "")
            for a in soup.select("div.mw-parser-output a[href^='/wiki/']")
            if ":" not in a["href"]]   # exclui namespaces especiais
G = nx.DiGraph()
visitados = set()
fila = ["Python_(linguagem_de_programação)"]
\end{code}
\end{frame}

\begin{frame}[c, fragile]{PageRank — implementação iterativa}

\begin{code}[NetworkX: PageRank em grafo da Wikipedia]{python}
while fila and len(visitados) < 30:   # limitar para demo
    pagina = fila.pop(0)
    if pagina in visitados: continue
    visitados.add(pagina)
    for dest in links_da_pagina(pagina)[:10]:
        G.add_edge(pagina, dest)
        fila.append(dest)
pr = nx.pagerank(G, alpha=0.85, max_iter=100)
top5 = sorted(pr, key=pr.get, reverse=True)[:5]
\end{code}
\end{frame}


\begin{frame}[c, fragile]{PageRank — do zero, sem bibliotecas}

\begin{code}[Iteração de potências — entendendo o algoritmo]{python}
    # Teleporte: internauta pula para página aleatória com prob. (1-alpha)
    r = np.ones(N) / N
    for _ in range(max_iter):
        r_novo = alpha * M @ r + (1 - alpha) / N
        if np.linalg.norm(r_novo - r, 1) < tol:
            break
        r = r_novo
    return dict(zip(nós, r))
\end{code}
\end{frame}

\begin{frame}[c, fragile]{PageRank — do zero, sem bibliotecas}

\begin{code}[Iteração de potências — entendendo o algoritmo]{python}
    # Teleporte: internauta pula para página aleatória com prob. (1-alpha)
    r = np.ones(N) / N
    for _ in range(max_iter):
        r_novo = alpha * M @ r + (1 - alpha) / N
        if np.linalg.norm(r_novo - r, 1) < tol:
            break
        r = r_novo
    return dict(zip(nós, r))
\end{code}
\end{frame}



% ============================================================
\section{LDA — Modelagem de Tópicos}
% ============================================================

\begin{frame}[c]{LDA — Latent Dirichlet Allocation}

\begin{columns}
\column{0.5\textwidth}
\textbf{\color{nes_dark_purple}A intuição}
\begin{itemize}
    \item Um documento é uma \highlight{mistura de tópicos}
    \item Cada tópico é uma \highlight{distribuição sobre palavras}
    \item Exemplo: documento sobre esporte pode ser 70\% tópico ``futebol'', 30\% tópico ``saúde''
    \item LDA \emph{infere} os tópicos latentes a partir das co-ocorrências de palavras
\end{itemize}



\column{0.5\textwidth}
\textbf{\color{nes_dark_purple}Processo gerador (modelo)}
\begin{enumerate}
    \item Para cada tópico $k$: sorteia distribuição sobre palavras $\phi_k \sim \text{Dir}(\beta)$
    \item Para cada documento $d$: sorteia mistura de tópicos $\theta_d \sim \text{Dir}(\alpha)$
    \item Para cada palavra $w$ em $d$: sorteia tópico $z \sim \theta_d$, depois palavra $w \sim \phi_z$
\end{enumerate}
\end{columns}
\end{frame}

\begin{frame}[c]{LDA — Latent Dirichlet Allocation}

\textbf{\color{nes_dark_purple}O que Dirichlet tem a ver?}
\begin{itemize}
    \item Distribuição de Dirichlet é uma \highlight{distribuição sobre distribuições}
    \item Parâmetro $\alpha$ pequeno $\Rightarrow$ documentos com poucos tópicos dominantes
    \item Parâmetro $\alpha$ grande $\Rightarrow$ mistura uniforme de tópicos
    \item Análogo para $\beta$ e as palavras nos tópicos
\end{itemize}
\vspace{0.4em}
\begin{display}[Inferência]
Observamos as palavras; queremos inferir $\theta$ e $\phi$. Usamos Gibbs Sampling ou Variational Bayes.
\end{display}
\end{frame}

\begin{frame}[c]{LDA — processo gerador}

\begin{figure}
    \includegraphics[height=.78\textheight]{lda_plate.jpg}
\end{figure}

\end{frame}


\begin{frame}[c, fragile]{LDA na prática}

\begin{code}[scikit-learn: LDA em corpus de notícias]{python}
from sklearn.decomposition import LatentDirichletAllocation
from sklearn.feature_extraction.text import CountVectorizer
=
# LDA usa contagens (não TF-IDF)
vec = CountVectorizer(max_df=0.95, min_df=2,
                      max_features=1000, stop_words="english")
X_counts = vec.fit_transform(corpus)
lda = LatentDirichletAllocation(n_components=5,     # número de tópicos
                                max_iter=20,
                                random_state=42)
lda.fit(X_counts)
# Palavras mais relevantes por tópico
termos = vec.get_feature_names_out()
for t, componente in enumerate(lda.components_):
    top_idx = componente.argsort()[::-1][:8]
    print(f"Tópico {t}: {' | '.join(termos[top_idx])}")
doc_topicos = lda.transform(X_counts)   # shape: (n_docs, n_topics)
\end{code}

\end{frame}


% ============================================================
\section{Busca semântica}
% ============================================================

\begin{frame}[c]{De TF-IDF a embeddings — busca semântica}

\begin{columns}
\column{0.6\textwidth}
\textbf{\color{nes_dark_purple}Limitações do TF-IDF}
\begin{itemize}
    \item Baseado em \highlight{correspondência exata} de tokens
    \item ``automóvel'' e ``carro'' são completamente distintos
    \item Sem noção de contexto ou sinônimos
    \item Funciona bem para recuperação léxica; mal para semântica
\end{itemize}

\vspace{0.4em}
\textbf{\color{nes_dark_purple}Sentence Transformers}
\begin{itemize}
    \item Modelo BERT ajustado para produzir \highlight{embeddings de sentenças}
    \item Documentos similares ficam próximos no espaço vetorial
    \item Busca = encontrar vizinhos mais próximos no espaço de embeddings
\end{itemize}

\column{0.4\textwidth}
\textbf{\color{nes_dark_purple}Pipeline completo}
\begin{enumerate}
    \item Gerar embeddings dos documentos com \texttt{sentence-transformers}
    \item Armazenar no banco (PostgreSQL com \texttt{pgvector} ou SQLite com \texttt{sqlite-vec})
    \item Na busca: gerar embedding da query e fazer busca por vizinhos mais próximos (\highlight{ANN})
    \item Retornar documentos ordenados por similaridade de cosseno
\end{enumerate}
\end{columns}

\end{frame}


\begin{frame}[c, fragile]{Busca semântica com sentence-transformers}

\begin{code}[Embeddings + SQLite: buscador semântico]{python}
from sentence_transformers import SentenceTransformer
import numpy as np, sqlite3, json

model = SentenceTransformer("paraphrase-multilingual-MiniLM-L12-v2")

# Indexar documentos
conn = sqlite3.connect("semantico.db")
conn.execute("CREATE TABLE IF NOT EXISTS docs "
             "(id INTEGER PRIMARY KEY, titulo TEXT, "
             " conteudo TEXT, embedding TEXT)")

for titulo, texto in documentos:
    emb = model.encode(texto).tolist()
    conn.execute("INSERT INTO docs (titulo, conteudo, embedding) VALUES (?,?,?)",
                 (titulo, texto, json.dumps(emb)))
conn.commit()
\end{code}

\end{frame}

\begin{frame}[c, fragile]{Busca semântica com sentence-transformers}

\begin{code}[Embeddings + SQLite: buscador semântico]{python}
# Buscar
def buscar_semantico(query, top_k=5):
    q_emb = model.encode(query)
    rows  = conn.execute("SELECT id, titulo, embedding FROM docs").fetchall()
    scores = [(r[0], r[1],
               float(np.dot(q_emb, json.loads(r[2])) /
               (np.linalg.norm(q_emb) * np.linalg.norm(json.loads(r[2])))))
              for r in rows]
    return sorted(scores, key=lambda x: -x[2])[:top_k]
\end{code}

\end{frame}

\begin{frame}[c, noframenumbering, plain]
    \frametitle{~}
        \vfill
        \begin{center}
            {\Huge Obrigado!}\vspace{1.5em}\\
            {\Large \highlight{Alguma dúvida?}}\\
        \end{center}
        \vfill
        \begin{center}
            {\small Agora vamos para os exercícios!}
        \end{center}
\end{frame}
\end{document}